% ============================================================
\chapter{Bitopological Spaces}
\label{ch:bitopological}
% ============================================================

In 1963, John C.\ Kelly proposed a bold idea: equip the same set with \emph{two} distinct topologies and study their interaction. The idea is far from gratuitous---it arises naturally from quasi-metrics, where the ``forward'' and ``backward'' topologies coexist, and from order theory, where upper and lower topologies on an ordered set capture complementary aspects. Nachbin's duality between compact ordered spaces and stably compact spaces shows that this bitopological structure is not an artefact, but a deep phenomenon.

\begin{intuition}
A bitopological space carries \emph{two} topologies on the same set, capturing a natural
duality. This structure arises naturally from quasi-metrics (forward and conjugate
topologies), from order theory (upper and lower topologies), and in Nachbin's duality
between compact ordered spaces and stably compact spaces.
\end{intuition}

% ------------------------------------------------------------
\section{Fundamental Definitions}
% ------------------------------------------------------------

\begin{definition}[Bitopological space]
A \emph{bitopological space} is a triple $(X, \tau_1, \tau_2)$ where $X$ is a set and
$\tau_1, \tau_2$ are two topologies on $X$. We call $\tau_1$ the \emph{primary topology}
and $\tau_2$ the \emph{secondary} (or \emph{dual}) topology.
\end{definition}

\begin{definition}[Bicontinuous map]
Let $(X, \tau_1, \tau_2)$ and $(Y, \sigma_1, \sigma_2)$ be bitopological spaces. A map
$f\colon X \to Y$ is \emph{bicontinuous} if $f$ is continuous from $(X,\tau_1)$ to
$(Y,\sigma_1)$ and from $(X,\tau_2)$ to $(Y,\sigma_2)$.
\end{definition}

\begin{definition}[Category $\mathbf{BiTop}$]
The category $\mathbf{BiTop}$ has bitopological spaces as objects and bicontinuous maps as
morphisms. A \emph{bi-homeomorphism} is an isomorphism in $\mathbf{BiTop}$.
\end{definition}

% ------------------------------------------------------------
\section{Fundamental Examples}
% ------------------------------------------------------------

\begin{example}[Bitopological space of a quasi-metric]
Let $(X,d)$ be a quasi-metric space. The associated bitopological space is
$(X, \tau_d, \tau_{d^{-1}})$, where $\tau_d$ is the forward topology and
$\tau_{d^{-1}}$ is the conjugate topology.
\end{example}

\begin{example}[Ordered spaces]
Let $(X,\leq)$ be a partially ordered set. The Alexandrov bitopological space is
$(X, \tau_{\mathrm{up}}, \tau_{\mathrm{down}})$, where $\tau_{\mathrm{up}}$ is the
Alexandrov topology (opens = upward-closed sets) and $\tau_{\mathrm{down}}$ is the dual
topology (opens = downward-closed sets).
\end{example}

\begin{example}[Ordered real line]
On $\R$, the topology of rightward open sets $\{(a,+\infty) : a \in \R\} \cup
\{\emptyset, \R\}$ and the topology of leftward open sets
$\{(-\infty, b) : b \in \R\} \cup \{\emptyset, \R\}$ define a bitopological space. Their
join is the usual topology of $\R$.
\end{example}

% ------------------------------------------------------------
\section{Bitopological Separation Axioms}
% ------------------------------------------------------------

\begin{definition}[Kelly separation]
A bitopological space $(X, \tau_1, \tau_2)$ is:
\begin{itemize}
  \item \emph{Pairwise $T_0$} if for all $x \neq y$, there exists an open in $\tau_1$ or
    $\tau_2$ containing one but not the other.
  \item \emph{Pairwise $T_1$} if for all $x \neq y$, there exists a $\tau_1$-open
    containing $x$ but not $y$, and a $\tau_2$-open containing $y$ but not $x$ (or
    vice versa).
  \item \emph{Pairwise $T_2$} (or \emph{pairwise Hausdorff}) if for all $x \neq y$,
    there exist a $\tau_1$-open $U \ni x$ and a $\tau_2$-open $V \ni y$ with
    $U \cap V = \emptyset$.
\end{itemize}
\end{definition}

\begin{theorem}[Pairwise $T_2$ and quasi-metrics]
If $(X,d)$ is a $T_0$-separated quasi-metric, then $(X,\tau_d, \tau_{d^{-1}})$ is
pairwise $T_2$.
\end{theorem}

\begin{proof}
Let $x \neq y$. Then $d(x,y) > 0$ or $d(y,x) > 0$. Suppose $d(x,y) > 0$ and set
$\varepsilon = d(x,y)/2 > 0$. Then $U = B_d(x,\varepsilon)$ is a $\tau_d$-open containing
$x$ and $V = B_{d^{-1}}(y,\varepsilon)$ is a $\tau_{d^{-1}}$-open containing $y$. If
$z \in U \cap V$, then $d(x,z) < \varepsilon$ and $d^{-1}(y,z) = d(z,y) < \varepsilon$,
so $d(x,y) \leq d(x,z) + d(z,y) < 2\varepsilon = d(x,y)$, a contradiction. Hence
$U \cap V = \emptyset$.
\end{proof}

% ------------------------------------------------------------
\section{Pairwise Normality}
% ------------------------------------------------------------

\begin{definition}[Pairwise normality]
A bitopological space $(X,\tau_1,\tau_2)$ is \emph{pairwise normal} if for every
$\tau_1$-closed $F_1$ and $\tau_2$-closed $F_2$ with $F_1 \cap F_2 = \emptyset$, there
exist a $\tau_2$-open $U_2 \supseteq F_1$ and a $\tau_1$-open $U_1 \supseteq F_2$ with
$U_1 \cap U_2 = \emptyset$.
\end{definition}

\begin{theorem}[Bitopological Urysohn lemma]
If $(X,\tau_1,\tau_2)$ is pairwise normal and pairwise $T_1$, then for every disjoint pair
of a $\tau_1$-closed $F_1$ and a $\tau_2$-closed $F_2$, there exists a function
$f\colon X \to [0,1]$ such that:
\begin{itemize}
  \item $f \equiv 0$ on $F_1$ and $f \equiv 1$ on $F_2$;
  \item $f$ is lower semicontinuous for $\tau_1$ and upper semicontinuous for $\tau_2$.
\end{itemize}
\end{theorem}

\begin{proof}
The proof follows the classical Urysohn scheme. By induction on dyadic rationals
$r \in [0,1] \cap \Z[1/2]$, construct sets $(A_r)$ such that:
\begin{itemize}
  \item $F_1 \subseteq A_0$ and $A_1 = X \setminus F_2$;
  \item $A_r$ is $\tau_2$-open for $r < 1$;
  \item $r < s \Rightarrow \overline{A_r}^{\tau_1} \subseteq A_s$.
\end{itemize}
Existence follows from pairwise normality. Define $f(x) = \inf\{r : x \in A_r\}$ and
verify the semicontinuity properties.
\end{proof}

% ------------------------------------------------------------
\section{Bitopological Compactness}
% ------------------------------------------------------------

\begin{definition}[Joint compactness]
A bitopological space $(X,\tau_1,\tau_2)$ is \emph{jointly compact} if the space
$(X, \tau_1 \vee \tau_2)$ is compact.
\end{definition}

\begin{definition}[Pairwise compactness]
A bitopological space is \emph{pairwise compact} if every cover consisting of opens from
$\tau_1$ and $\tau_2$ has a finite subcover.
\end{definition}

\begin{proposition}
Pairwise compactness implies joint compactness. The converse holds when $\tau_1$ and
$\tau_2$ form a ``compatible pair.''
\end{proposition}

% ------------------------------------------------------------
\section{Nachbin Duality}
% ------------------------------------------------------------

\begin{definition}[Compact ordered space]
A \emph{compact ordered space} (in the sense of Nachbin) is a triple $(X, \tau, \leq)$
where:
\begin{enumerate}
  \item $(X,\tau)$ is a compact Hausdorff space;
  \item $\leq$ is a partial order that is closed in $X \times X$ (for the product
    topology).
\end{enumerate}
\end{definition}

\begin{theorem}[Nachbin duality]
There is an equivalence of categories between:
\begin{enumerate}
  \item the category of compact ordered spaces (with continuous monotone maps);
  \item the category of \emph{pairwise compact pairwise $T_2$ jointly $T_1$}
    bitopological spaces (with bicontinuous maps).
\end{enumerate}
\end{theorem}

\begin{proof}[Proof sketch]
$(\Rightarrow)$ Given a compact ordered space $(X,\tau,\leq)$, define:
\begin{itemize}
  \item $\tau_1 =$ the topology of upward-closed opens (upper topology within $\tau$);
  \item $\tau_2 =$ the topology of downward-closed opens.
\end{itemize}
More precisely, $\tau_1$ is generated by opens $U \in \tau$ with $U = \uparrow\! U$ and
$\tau_2$ by opens $V \in \tau$ with $V = \downarrow\! V$. Nachbin's theorem gives
$\tau_1 \vee \tau_2 = \tau$.

$(\Leftarrow)$ Given a bitopological space $(X,\tau_1,\tau_2)$ of the described type, set
$\tau = \tau_1 \vee \tau_2$ and define $x \leq y \Leftrightarrow
x \in \overline{\{y\}}^{\tau_1}$.
\end{proof}

% ------------------------------------------------------------
\section{Connection with Stably Compact Spaces}
% ------------------------------------------------------------

\begin{theorem}[Stably compact and bitopological spaces]
Let $X$ be a stably compact space. Then:
\begin{enumerate}
  \item The co-compact topology $\tau^d$ (generated by complements of compact saturated
    sets) determines a bitopological space $(X, \tau, \tau^d)$.
  \item The triple $(X, \tau \vee \tau^d, \leq)$ is a compact ordered space.
  \item This construction establishes an equivalence between the category of stably compact
    spaces and a subcategory of compact ordered spaces.
\end{enumerate}
\end{theorem}

\begin{remark}
This is one of the deepest results in asymmetric topology: stably compact spaces, compact
ordered spaces, and certain bitopological spaces are three facets of the same phenomenon.
\end{remark}

% ------------------------------------------------------------
\section{Functors Between $\mathbf{BiTop}$ and $\mathbf{Top}$}
% ------------------------------------------------------------

\begin{proposition}[Forgetful and symmetrization functors]
The following functors are available:
\begin{enumerate}
  \item $U_1, U_2\colon \mathbf{BiTop} \to \mathbf{Top}$, projection onto each topology.
  \item $J\colon \mathbf{BiTop} \to \mathbf{Top}$, sending $(X,\tau_1,\tau_2)$ to
    $(X, \tau_1 \vee \tau_2)$ (join topology).
  \item $\Delta\colon \mathbf{Top} \to \mathbf{BiTop}$, sending $(X,\tau)$ to
    $(X,\tau,\tau)$ (diagonal doubling).
\end{enumerate}
The functor $\Delta$ is left and right adjoint to $J$ in certain subcategories.
\end{proposition}

% ------------------------------------------------------------
\section{Bitopological Spaces and Lattices}
% ------------------------------------------------------------

\begin{definition}[Biframe]
A \emph{biframe} is a triple $(L, L_1, L_2)$ where $L$ is a frame and $L_1, L_2$ are
subframes of $L$ that generate $L$: every element of $L$ is a join of finite meets of
elements from $L_1 \cup L_2$.
\end{definition}

\begin{remark}
Biframes are the ``pointless'' analogue of bitopological spaces, just as frames are the
analogue of topological spaces. The theory of bilocales developed by Banaschewski,
Br\"ummer, and Hardie provides a rich algebraic framework for bitopology.
\end{remark}

% ------------------------------------------------------------
\section{Exercises}
% ------------------------------------------------------------

\begin{exercise}
Show that for any quasi-metric $d$, the bitopological space $(X,\tau_d,\tau_{d^{-1}})$ is
pairwise $T_0$ if and only if $d$ is $T_0$-separated.
\end{exercise}

\begin{exercise}
Let $(X,\tau_1,\tau_2)$ be pairwise $T_2$. Show that the join topology $\tau_1 \vee \tau_2$
is Hausdorff.
\end{exercise}

\begin{exercise}
Show that every compact ordered space is pairwise normal.
\end{exercise}

\begin{exercise}
Let $X = [0,1]$ with the usual order. Describe explicitly the topologies
$\tau_{\mathrm{up}}$, $\tau_{\mathrm{down}}$, the join topology, and verify that
$(X, \tau_{\mathrm{up}} \vee \tau_{\mathrm{down}}, \leq)$ is a compact ordered space.
\end{exercise}

\begin{exercise}
Show that the category $\mathbf{BiTop}$ is complete and cocomplete. Describe products and
coproducts.
\end{exercise}

\begin{exercise}[Priestley duality]
Let $D$ be a bounded distributive lattice and $X$ the associated Priestley space (compact
ordered totally order-separated space). Describe the corresponding bitopological space via
Nachbin duality and show that one recovers the spectral space of Stone duality.
\end{exercise}
